\documentclass[12pt,a4paper]{article}
\usepackage[utf8]{inputenc}
\usepackage[spanish]{babel}
\usepackage{geometry}
\usepackage{graphicx}
\usepackage{hyperref}
\usepackage{enumitem}
\usepackage{booktabs}
\usepackage{listings}
\usepackage{xcolor}
\usepackage{longtable}
\usepackage{pdflscape}

\geometry{margin=2.5cm}
\hypersetup{
    colorlinks=true,
    linkcolor=blue,
    urlcolor=blue,
    citecolor=blue
}

\lstset{
    basicstyle=\ttfamily\small,
    breaklines=true,
    frame=single,
    numbers=left,
    numberstyle=\tiny,
    showstringspaces=false
}

\title{\textbf{Documento de Diseño Técnico}\\
\large Secure Fair: Sistema de Registro para Ferias de Servicio Social}
\author{Equipo de Desarrollo}
\date{Febrero 2026}

\begin{document}

\maketitle
\newpage

\tableofcontents
\newpage

\section{Arquitectura del Sistema}

\subsection{Arquitectura General}

El sistema sigue una arquitectura de tres capas:

\begin{verbatim}
+------------------------------------------------------+
|                   CAPA DE PRESENTACIÓN                |
|                                                        |
|  +-------------+  +-------------+  +-------------+  |
|  |  Student    |  |   Socio     |  |   Admin     |  |
|  |  React App  |  |  React App  |  |  React App  |  |
|  +-------------+  +-------------+  +-------------+  |
|                                                        |
|              React + TypeScript + Vite                |
+--------------------+---------------------------------+
                     | HTTPS / REST API
                     |
+--------------------v---------------------------------+
|                   CAPA DE LÓGICA                      |
|                                                        |
|              FastAPI + Python 3.11+                   |
|                                                        |
|  +----------------------------------------------+   |
|  |           API Endpoints                       |   |
|  |  /auth  /admin  /student  /socio             |   |
|  +----------------------------------------------+   |
|                                                        |
|  +----------------------------------------------+   |
|  |       Business Logic & Services               |   |
|  |  - Enrollment Service                         |   |
|  |  - Code Generation Service                    |   |
|  |  - Crypto Service (Ed25519)                   |   |
|  |  - Check-in Service                           |   |
|  +----------------------------------------------+   |
|                                                        |
|  +----------------------------------------------+   |
|  |          ORM & Data Access                    |   |
|  |          SQLAlchemy 2.0                       |   |
|  +----------------------------------------------+   |
+--------------------+---------------------------------+
                     | SQL / Transacciones
                     |
+--------------------v---------------------------------+
|                   CAPA DE DATOS                       |
|                                                        |
|                  PostgreSQL 15+                       |
|                                                        |
|  - Tablas normalizadas                                |
|  - Restricciones de integridad                        |
|  - Índices optimizados                                |
|  - Transacciones ACID                                 |
+------------------------------------------------------+
\end{verbatim}

\subsection{Diagrama de Componentes}

\begin{verbatim}
Frontend Components:
+-- Auth Module
|   +-- Login
|   +-- Role Guard
|   +-- Token Manager
+-- Student Module
|   +-- Slot Registration
|   +-- QR Display
|   +-- Code Redemption
|   +-- Enrollment Status
+-- Socio Module
|   +-- Project Dashboard
|   +-- Code Generator
|   +-- Enrollment List
|   +-- Export Tools
+-- Admin Module
    +-- CRUD Organizations
    +-- CRUD Projects
    +-- CRUD Slots
    +-- Check-in Interface
    +-- Analytics Dashboard
    +-- Master Export

Backend Services:
+-- Auth Service
|   +-- JWT Generation
|   +-- Password Hashing (Argon2)
|   +-- Role Verification
+-- Enrollment Service
|   +-- Capacity Check
|   +-- Duplicate Prevention
|   +-- Transaction Management
+-- Code Service
|   +-- Code Generation
|   +-- HMAC Hashing
|   +-- Expiry Management
+-- Crypto Service
|   +-- Ed25519 Key Management
|   +-- Receipt Signing
|   +-- Signature Verification
+-- Export Service
    +-- CSV Generation
    +-- XLSX Generation
\end{verbatim}

\subsection{Modelo de Autenticación}

\subsubsection{Flujo de Autenticación}

\begin{enumerate}
    \item Usuario envía credenciales a POST /auth/login
    \item Backend verifica credenciales con hash Argon2
    \item Si válido, genera JWT con payload:
    \begin{verbatim}
    {
        "sub": "<user_id>",
        "email": "<email>",
        "role": "STUDENT|SOCIO|ADMIN",
        "exp": <timestamp>,
        "iat": <timestamp>
    }
    \end{verbatim}
    \item Cliente almacena JWT en localStorage
    \item Cada request incluye header: \texttt{Authorization: Bearer <token>}
    \item Backend valida firma y expiración en cada endpoint protegido
\end{enumerate}

\subsubsection{Estrategia de Tokens}

\begin{itemize}[leftmargin=*]
    \item \textbf{Access Token:} Duración 15 minutos
    \item \textbf{Algoritmo:} HS256 (HMAC con SHA-256)
    \item \textbf{Secret Key:} Variable de entorno (mínimo 256 bits)
    \item \textbf{Renovación:} Re-login requerido tras expiración (simplificado para MVP)
\end{itemize}

\subsection{Diagrama de Despliegue}

\begin{verbatim}
                        Internet
                           |
                           v
                   +---------------+
                   |  CDN / Nginx  |
                   |   (Frontend)  |
                   +-------+-------+
                           |
        +------------------+------------------+
        |                  |                  |
        v                  v                  v
+---------------+  +---------------+  +--------------+
|   Student     |  |     Socio     |  |    Admin     |
|   SPA         |  |     SPA       |  |    SPA       |
| (Vercel/Netlify)| | (Vercel/Netlify)| |(Vercel/Netlify)|
+-------+-------+  +-------+-------+  +------+-------+
        |                  |                  |
        +------------------+------------------+
                           | HTTPS/TLS 1.3
                           v
                   +---------------+
                   |  FastAPI App  |
                   | (Render/Fly.io)|
                   |  + Gunicorn   |
                   +-------+-------+
                           | Internal Network
                           v
                   +---------------+
                   |  PostgreSQL   |
                   |   Database    |
                   | (Managed DB)  |
                   +---------------+
\end{verbatim}

\section{Esquema de Base de Datos}

\subsection{Diagrama de Entidad-Relación}

\begin{verbatim}
+-------------+         +--------------+         +--------------+
|    users    |         |organizations |         |fair_periods  |
+-------------+         +--------------+         +--------------+
| id PK       |         | id PK        |         | id PK        |
| email UQ    |         | name         |         | name         |
| password_h..|         | created_at   |         | starts_at    |
| role        |         +------+-------+         | ends_at      |
| is_active   |                |                 | is_active    |
| created_at  |                |                 +------+-------+
+------+------+                |                        |
       |                       |                        |
       | 1                     |                        |
       |                       | N                      | 1
       |                       v                        |
       |              +--------------+                 |
       |              |   projects   |<----------------+
       |              +--------------+         N
       |              | id PK        |
       |              | org_id FK    |
       |              | period_id FK |
       |              | name         |
       |              | description  |
       |              | rules_text   |
       |              | capacity     |
       |              | is_active    |
       |              +------+-------+
       |                     |
       |                     | N
       |                     |
       |              +------v--------------+
       |              |project_socio_users  |
       |              +---------------------+
       |     +--------+ project_id FK  PK   |
       |     |        | user_id FK  PK      |
       |     |        +---------------------+
       |     |
       | 1   | N              +--------------+
       v     +--------------->|  time_slots  |
+-------------+               +--------------+
|  students   |               | id PK        |
+-------------+               | period_id FK |
| user_id PK,F|               | starts_at    |
| matricula UQ|               | ends_at      |
| full_name   |               | capacity     |
| phone       |               | location     |
+------+------+               +------+-------+
       |                             |
       | 1                           | 1
       |                             |
       |                             | N
       |                      +------v-------------+
       |                      |slot_registrations  |
       |                      +--------------------+
       |                      | id PK              |
       |            +---------+ slot_id FK         |
       |            |         | student_user_id FK |
       |            |         | created_at         |
       |            |         +--------+-----------+
       |            |                  |
       |            |                  | 1
       |            |                  |
       |            |                  | 1
       |            |           +------v------+
       |            |           |  checkins   |
       |            |           +-------------+
       |            |           | id PK       |
       |            |           | slot_reg.FK |
       |            |           | checked_in.a|
       |            |           | checked_by..|
       |            |           +-------------+
       |            |
       | 1          |
       |            |
       | N          |
+------v---------+  |          +--------------+
|  enrollments   |  |          |project_codes |
+----------------+  |          +--------------+
| id PK          |  |          | id PK        |
| period_id FK   |  |          | project_id FK|
| project_id FK  |<-+          | code_hash    |
| student_u.. FK |             | expires_at   |
| accepted_ru..  |             | issued_by FK |
| receipt_sig.   |             | used_by_st.. |
| created_at     |             | used_at      |
+----------------+             +--------------+
UQ: (student_user_id, period_id)
\end{verbatim}

\subsection{Definiciones de Tablas}

\subsubsection{users}

\begin{lstlisting}[language=SQL]
CREATE TABLE users (
    id UUID PRIMARY KEY DEFAULT gen_random_uuid(),
    email VARCHAR(255) NOT NULL UNIQUE,
    password_hash TEXT NOT NULL,
    role VARCHAR(20) NOT NULL CHECK (role IN ('ADMIN', 'SOCIO', 'STUDENT')),
    is_active BOOLEAN DEFAULT TRUE,
    created_at TIMESTAMP DEFAULT CURRENT_TIMESTAMP
);

CREATE INDEX idx_users_email ON users(email);
CREATE INDEX idx_users_role ON users(role);
\end{lstlisting}

\subsubsection{organizations}

\begin{lstlisting}[language=SQL]
CREATE TABLE organizations (
    id UUID PRIMARY KEY DEFAULT gen_random_uuid(),
    name VARCHAR(255) NOT NULL,
    created_at TIMESTAMP DEFAULT CURRENT_TIMESTAMP
);
\end{lstlisting}

\subsubsection{fair\_periods}

\begin{lstlisting}[language=SQL]
CREATE TABLE fair_periods (
    id UUID PRIMARY KEY DEFAULT gen_random_uuid(),
    name VARCHAR(255) NOT NULL,
    starts_at TIMESTAMP NOT NULL,
    ends_at TIMESTAMP NOT NULL,
    is_active BOOLEAN DEFAULT TRUE,
    CHECK (ends_at > starts_at)
);

CREATE INDEX idx_fair_periods_active ON fair_periods(is_active);
\end{lstlisting}

\subsubsection{projects}

\begin{lstlisting}[language=SQL]
CREATE TABLE projects (
    id UUID PRIMARY KEY DEFAULT gen_random_uuid(),
    organization_id UUID NOT NULL REFERENCES organizations(id) ON DELETE CASCADE,
    period_id UUID NOT NULL REFERENCES fair_periods(id) ON DELETE CASCADE,
    name VARCHAR(255) NOT NULL,
    description TEXT,
    rules_text TEXT,
    capacity INTEGER CHECK (capacity IS NULL OR capacity > 0),
    is_active BOOLEAN DEFAULT TRUE,
    created_at TIMESTAMP DEFAULT CURRENT_TIMESTAMP
);

CREATE INDEX idx_projects_org ON projects(organization_id);
CREATE INDEX idx_projects_period ON projects(period_id);
CREATE INDEX idx_projects_active ON projects(is_active);
\end{lstlisting}

\subsubsection{project\_socio\_users}

\begin{lstlisting}[language=SQL]
CREATE TABLE project_socio_users (
    project_id UUID NOT NULL REFERENCES projects(id) ON DELETE CASCADE,
    user_id UUID NOT NULL REFERENCES users(id) ON DELETE CASCADE,
    PRIMARY KEY (project_id, user_id)
);

CREATE INDEX idx_proj_socio_user ON project_socio_users(user_id);
\end{lstlisting}

\subsubsection{time\_slots}

\begin{lstlisting}[language=SQL]
CREATE TABLE time_slots (
    id UUID PRIMARY KEY DEFAULT gen_random_uuid(),
    period_id UUID NOT NULL REFERENCES fair_periods(id) ON DELETE CASCADE,
    starts_at TIMESTAMP NOT NULL,
    ends_at TIMESTAMP NOT NULL,
    capacity INTEGER NOT NULL CHECK (capacity > 0),
    location VARCHAR(255),
    CHECK (ends_at > starts_at)
);

CREATE INDEX idx_slots_period ON time_slots(period_id, starts_at);
\end{lstlisting}

\subsubsection{students}

\begin{lstlisting}[language=SQL]
CREATE TABLE students (
    user_id UUID PRIMARY KEY REFERENCES users(id) ON DELETE CASCADE,
    matricula VARCHAR(50) NOT NULL UNIQUE,
    full_name VARCHAR(255) NOT NULL,
    phone VARCHAR(20),
    non_tec_email VARCHAR(255)
);

CREATE INDEX idx_students_matricula ON students(matricula);
\end{lstlisting}

\subsubsection{slot\_registrations}

\begin{lstlisting}[language=SQL]
CREATE TABLE slot_registrations (
    id UUID PRIMARY KEY DEFAULT gen_random_uuid(),
    slot_id UUID NOT NULL REFERENCES time_slots(id) ON DELETE CASCADE,
    student_user_id UUID NOT NULL REFERENCES users(id) ON DELETE CASCADE,
    created_at TIMESTAMP DEFAULT CURRENT_TIMESTAMP,
    UNIQUE (slot_id, student_user_id)
);

CREATE INDEX idx_slot_reg_student ON slot_registrations(student_user_id);
CREATE INDEX idx_slot_reg_slot ON slot_registrations(slot_id);
\end{lstlisting}

\subsubsection{checkins}

\begin{lstlisting}[language=SQL]
CREATE TABLE checkins (
    id UUID PRIMARY KEY DEFAULT gen_random_uuid(),
    slot_registration_id UUID NOT NULL UNIQUE 
        REFERENCES slot_registrations(id) ON DELETE CASCADE,
    checked_in_at TIMESTAMP DEFAULT CURRENT_TIMESTAMP,
    checked_in_by_admin_user_id UUID REFERENCES users(id)
);

CREATE INDEX idx_checkins_slot_reg ON checkins(slot_registration_id);
\end{lstlisting}

\subsubsection{enrollments}

\begin{lstlisting}[language=SQL]
CREATE TABLE enrollments (
    id UUID PRIMARY KEY DEFAULT gen_random_uuid(),
    period_id UUID NOT NULL REFERENCES fair_periods(id) ON DELETE CASCADE,
    project_id UUID NOT NULL REFERENCES projects(id) ON DELETE CASCADE,
    student_user_id UUID NOT NULL REFERENCES users(id) ON DELETE CASCADE,
    accepted_rules_at TIMESTAMP DEFAULT CURRENT_TIMESTAMP,
    receipt_signature TEXT,
    created_at TIMESTAMP DEFAULT CURRENT_TIMESTAMP,
    UNIQUE (student_user_id, period_id)
);

CREATE INDEX idx_enrollments_project ON enrollments(project_id);
CREATE INDEX idx_enrollments_student ON enrollments(student_user_id);
CREATE INDEX idx_enrollments_period ON enrollments(period_id);
\end{lstlisting}

\subsubsection{project\_codes}

\begin{lstlisting}[language=SQL]
CREATE TABLE project_codes (
    id UUID PRIMARY KEY DEFAULT gen_random_uuid(),
    project_id UUID NOT NULL REFERENCES projects(id) ON DELETE CASCADE,
    code_hash TEXT NOT NULL,
    expires_at TIMESTAMP NOT NULL,
    issued_by_user_id UUID NOT NULL REFERENCES users(id),
    used_by_student_user_id UUID REFERENCES users(id),
    used_at TIMESTAMP,
    created_at TIMESTAMP DEFAULT CURRENT_TIMESTAMP
);

CREATE INDEX idx_proj_codes_project ON project_codes(project_id);
CREATE INDEX idx_proj_codes_expires ON project_codes(expires_at);
CREATE INDEX idx_proj_codes_used ON project_codes(used_at);
\end{lstlisting}

\subsection{Restricciones e Integridad}

\subsubsection{Restricciones Clave}

\begin{enumerate}
    \item \textbf{Inscripción única por periodo:} UNIQUE(student\_user\_id, period\_id) en enrollments
    \item \textbf{Registro único por slot:} UNIQUE(slot\_id, student\_user\_id) en slot\_registrations
    \item \textbf{Check-in único por registro:} UNIQUE(slot\_registration\_id) en checkins
    \item \textbf{Email único:} UNIQUE(email) en users
    \item \textbf{Matrícula única:} UNIQUE(matricula) en students
\end{enumerate}

\subsubsection{Validaciones CHECK}

\begin{itemize}[leftmargin=*]
    \item Roles válidos: CHECK (role IN ('ADMIN', 'SOCIO', 'STUDENT'))
    \item Capacidad positiva: CHECK (capacity > 0)
    \item Fechas coherentes: CHECK (ends\_at > starts\_at)
\end{itemize}

\subsection{Estrategia de Indexación}

\subsubsection{Índices Principales}

\begin{itemize}[leftmargin=*]
    \item \textbf{users.email:} Búsqueda rápida en login
    \item \textbf{students.matricula:} Búsqueda por matrícula
    \item \textbf{projects (organization\_id, period\_id):} Listados filtrados
    \item \textbf{time\_slots (period\_id, starts\_at):} Búsqueda de slots disponibles
    \item \textbf{enrollments (project\_id):} Conteo de inscritos por proyecto
    \item \textbf{project\_codes (expires\_at):} Limpieza de códigos expirados
\end{itemize}

\subsubsection{Índices Compuestos}

Para queries comunes como "obtener slots disponibles de un periodo ordenados por fecha", se recomienda:

\begin{lstlisting}[language=SQL]
CREATE INDEX idx_slots_period_time 
    ON time_slots(period_id, starts_at);
\end{lstlisting}

\section{Diseño de API}

\subsection{Convenciones Generales}

\subsubsection{Formato de Respuesta}

\begin{lstlisting}[language=json]
{
    "success": true,
    "data": { ... },
    "message": "Operation successful"
}
\end{lstlisting}

\subsubsection{Manejo de Errores}

\begin{lstlisting}[language=json]
{
    "success": false,
    "error": {
        "code": "ENROLLMENT_ALREADY_EXISTS",
        "message": "Ya estas inscrito en un proyecto para este periodo",
        "details": { ... }
    }
}
\end{lstlisting}

\subsubsection{Códigos de Estado HTTP}

\begin{itemize}[leftmargin=*]
    \item 200: Éxito en operación GET/PUT
    \item 201: Recurso creado exitosamente (POST)
    \item 400: Error de validación o regla de negocio
    \item 401: No autenticado
    \item 403: No autorizado (rol insuficiente)
    \item 404: Recurso no encontrado
    \item 409: Conflicto (ej. violación de unicidad)
    \item 500: Error interno del servidor
\end{itemize}

\subsection{Endpoints de Autenticación}

\subsubsection{POST /api/auth/login}

\textbf{Request:}
\begin{lstlisting}[language=json]
{
    "email": "student@example.com",
    "password": "SecurePass123"
}
\end{lstlisting}

\textbf{Response (200):}
\begin{lstlisting}[language=json]
{
    "success": true,
    "data": {
        "access_token": "eyJhbGciOiJIUzI1NiIsInR5cCI6IkpXVCJ9...",
        "token_type": "bearer",
        "user": {
            "id": "uuid",
            "email": "student@example.com",
            "role": "STUDENT"
        }
    }
}
\end{lstlisting}

\subsubsection{GET /api/auth/me}

\textbf{Headers:} Authorization: Bearer <token>

\textbf{Response (200):}
\begin{lstlisting}[language=json]
{
    "success": true,
    "data": {
        "id": "uuid",
        "email": "student@example.com",
        "role": "STUDENT",
        "is_active": true
    }
}
\end{lstlisting}

\subsection{Endpoints de Estudiante}

\subsubsection{GET /api/student/slots}

\textbf{Query Params:} period\_id (required)

\textbf{Response (200):}
\begin{lstlisting}[language=json]
{
    "success": true,
    "data": [
        {
            "id": "uuid",
            "starts_at": "2026-03-15T10:00:00Z",
            "ends_at": "2026-03-15T12:00:00Z",
            "capacity": 50,
            "registered_count": 42,
            "location": "Auditorio Principal",
            "available": true
        }
    ]
}
\end{lstlisting}

\subsubsection{POST /api/student/slot-registrations}

\textbf{Request:}
\begin{lstlisting}[language=json]
{
    "slot_id": "uuid"
}
\end{lstlisting}

\textbf{Response (201):}
\begin{lstlisting}[language=json]
{
    "success": true,
    "data": {
        "id": "uuid",
        "slot_id": "uuid",
        "created_at": "2026-03-10T15:30:00Z"
    },
    "message": "Registro exitoso"
}
\end{lstlisting}

\subsubsection{GET /api/student/slot-qr}

\textbf{Response (200):}
\begin{lstlisting}[language=json]
{
    "success": true,
    "data": {
        "qr_token": "eyJhbGciOiJIUzI1NiIsInR5cCI6IkpXVCJ9...",
        "qr_data_url": "data:image/png;base64,iVBORw0KG...",
        "slot_info": {
            "starts_at": "2026-03-15T10:00:00Z",
            "location": "Auditorio Principal"
        }
    }
}
\end{lstlisting}

\subsubsection{POST /api/student/enrollments/redeem}

\textbf{Request:}
\begin{lstlisting}[language=json]
{
    "code": "ABC123",
    "accept_rules": true
}
\end{lstlisting}

\textbf{Response (201):}
\begin{lstlisting}[language=json]
{
    "success": true,
    "data": {
        "enrollment_id": "uuid",
        "project_name": "Reforestacion Comunitaria",
        "organization_name": "EcoAmigos AC",
        "receipt": {
            "student_id": "uuid",
            "project_id": "uuid",
            "period_id": "uuid",
            "timestamp": "2026-03-15T10:45:00Z",
            "nonce": "random-nonce"
        },
        "signature": "base64-encoded-signature"
    },
    "message": "Inscripcion exitosa"
}
\end{lstlisting}

\textbf{Errores posibles:}
\begin{itemize}[leftmargin=*]
    \item 400: "Codigo invalido o expirado"
    \item 400: "No estas verificado (check-in). Primero presenta tu QR en la entrada"
    \item 409: "Ya estas inscrito en un proyecto para este periodo"
    \item 409: "Proyecto sin capacidad disponible"
\end{itemize}

\subsection{Endpoints de Socioformador}

\subsubsection{GET /api/socio/projects}

\textbf{Response (200):}
\begin{lstlisting}[language=json]
{
    "success": true,
    "data": [
        {
            "id": "uuid",
            "name": "Reforestacion Comunitaria",
            "organization_name": "EcoAmigos AC",
            "period_name": "Feb-Jun 2026",
            "capacity": 20,
            "enrolled_count": 15,
            "is_active": true
        }
    ]
}
\end{lstlisting}

\subsubsection{POST /api/socio/projects/\{project\_id\}/codes}

\textbf{Response (201):}
\begin{lstlisting}[language=json]
{
    "success": true,
    "data": {
        "code": "XK7M9P",
        "expires_at": "2026-03-15T10:47:00Z",
        "expires_in_seconds": 120
    },
    "message": "Muestra este codigo al estudiante"
}
\end{lstlisting}

\subsubsection{GET /api/socio/projects/\{project\_id\}/enrollments}

\textbf{Response (200):}
\begin{lstlisting}[language=json]
{
    "success": true,
    "data": [
        {
            "student_name": "Juan Perez",
            "matricula": "A01234567",
            "email": "juan@example.com",
            "phone": "5551234567",
            "enrolled_at": "2026-03-15T10:45:00Z"
        }
    ]
}
\end{lstlisting}

\subsubsection{GET /api/socio/projects/\{project\_id\}/enrollments/export}

\textbf{Response (200):} Archivo CSV/XLSX

\subsection{Endpoints de Administrador}

\subsubsection{Gestión de Periodos}

\textbf{GET /api/admin/periods}

\textbf{POST /api/admin/periods}

\textbf{Request:}
\begin{lstlisting}[language=json]
{
    "name": "Feb-Jun 2026",
    "starts_at": "2026-02-01T00:00:00Z",
    "ends_at": "2026-06-30T23:59:59Z",
    "is_active": true
}
\end{lstlisting}

\textbf{PUT /api/admin/periods/\{id\}}

\subsubsection{Gestión de Organizaciones}

\textbf{GET /api/admin/orgs}

\textbf{POST /api/admin/orgs}

\textbf{Request:}
\begin{lstlisting}[language=json]
{
    "name": "EcoAmigos AC"
}
\end{lstlisting}

\subsubsection{Gestión de Proyectos}

\textbf{GET /api/admin/projects?period\_id=uuid}

\textbf{POST /api/admin/projects}

\textbf{Request:}
\begin{lstlisting}[language=json]
{
    "organization_id": "uuid",
    "period_id": "uuid",
    "name": "Reforestacion Comunitaria",
    "description": "Proyecto de reforestacion...",
    "rules_text": "Compromiso de 480 horas...",
    "capacity": 20,
    "socio_user_ids": ["uuid1", "uuid2"]
}
\end{lstlisting}

\subsubsection{Gestión de Franjas Horarias}

\textbf{POST /api/admin/slots}

\textbf{Request:}
\begin{lstlisting}[language=json]
{
    "period_id": "uuid",
    "starts_at": "2026-03-15T10:00:00Z",
    "ends_at": "2026-03-15T12:00:00Z",
    "capacity": 50,
    "location": "Auditorio Principal"
}
\end{lstlisting}

\subsubsection{Check-in}

\textbf{POST /api/admin/checkin}

\textbf{Request:}
\begin{lstlisting}[language=json]
{
    "qr_token": "eyJhbGciOiJIUzI1NiIsInR5cCI6IkpXVCJ9..."
}
\end{lstlisting}

\textbf{Response (200):}
\begin{lstlisting}[language=json]
{
    "success": true,
    "data": {
        "checkin_id": "uuid",
        "student_name": "Juan Perez",
        "matricula": "A01234567",
        "slot_info": {
            "starts_at": "2026-03-15T10:00:00Z",
            "location": "Auditorio Principal"
        },
        "checked_in_at": "2026-03-15T09:55:00Z"
    },
    "message": "Check-in exitoso"
}
\end{lstlisting}

\subsubsection{Dashboard}

\textbf{GET /api/admin/dashboard?period\_id=uuid}

\textbf{Response (200):}
\begin{lstlisting}[language=json]
{
    "success": true,
    "data": {
        "total_slots": 10,
        "total_registrations": 450,
        "total_checkins": 420,
        "total_enrollments": 380,
        "slots_by_attendance": [
            {
                "slot_starts_at": "2026-03-15T10:00:00Z",
                "registered": 50,
                "checked_in": 48
            }
        ],
        "projects_by_enrollment": [
            {
                "project_name": "Reforestacion Comunitaria",
                "capacity": 20,
                "enrolled": 18,
                "fill_rate": 0.90
            }
        ]
    }
}
\end{lstlisting}

\subsubsection{Exportaciones}

\textbf{GET /api/admin/exports/master?period\_id=uuid}

\textbf{Response (200):} Archivo CSV/XLSX con columnas:
\begin{itemize}[leftmargin=*]
    \item Matricula, Nombre, Email, Telefono
    \item Slot (fecha/hora), Check-in (fecha/hora)
    \item Proyecto, Organizacion
    \item Fecha de inscripcion
\end{itemize}

\section{Reglas de Negocio Críticas}

\subsection{Lógica Transaccional de Inscripción}

\textbf{Pseudocódigo del endpoint POST /student/enrollments/redeem:}

\begin{lstlisting}[language=Python]
@router.post("/student/enrollments/redeem")
async def redeem_code(
    code: str,
    accept_rules: bool,
    current_user: User = Depends(get_current_student)
):
    if not accept_rules:
        raise HTTPException(400, "Debes aceptar las reglas")
    
    async with db.transaction():
        # 1. Verificar check-in
        checkin = await get_student_checkin(current_user.id)
        if not checkin:
            raise HTTPException(
                400, 
                "No estas verificado. Presenta tu QR en la entrada"
            )
        
        # 2. Buscar codigo valido
        code_hash = hmac_hash(code)
        code_record = await db.query(
            """
            SELECT * FROM project_codes
            WHERE code_hash = :hash
              AND expires_at > NOW()
              AND used_at IS NULL
            FOR UPDATE
            """
        ).first()
        
        if not code_record:
            raise HTTPException(400, "Codigo invalido o expirado")
        
        project = await get_project(code_record.project_id)
        
        # 3. Verificar no inscrito previamente
        existing = await db.query(
            """
            SELECT 1 FROM enrollments
            WHERE student_user_id = :student_id
              AND period_id = :period_id
            """
        ).first()
        
        if existing:
            raise HTTPException(
                409,
                "Ya estas inscrito en un proyecto para este periodo"
            )
        
        # 4. Verificar capacidad
        enrolled_count = await db.query(
            """
            SELECT COUNT(*) FROM enrollments
            WHERE project_id = :project_id
            FOR UPDATE
            """
        ).scalar()
        
        if project.capacity and enrolled_count >= project.capacity:
            raise HTTPException(409, "Proyecto sin capacidad")
        
        # 5. Crear inscripcion
        enrollment = await db.insert(
            enrollments,
            student_user_id=current_user.id,
            project_id=project.id,
            period_id=project.period_id,
            accepted_rules_at=datetime.utcnow()
        )
        
        # 6. Marcar codigo usado
        await db.execute(
            """
            UPDATE project_codes
            SET used_by_student_user_id = :student_id,
                used_at = NOW()
            WHERE id = :code_id
            """
        )
        
        # 7. Generar recibo firmado
        receipt = create_receipt(enrollment)
        signature = ed25519_sign(receipt)
        
        await db.execute(
            """
            UPDATE enrollments
            SET receipt_signature = :signature
            WHERE id = :enrollment_id
            """
        )
        
        # COMMIT automatico al salir del bloque transaction
    
    return {
        "enrollment_id": enrollment.id,
        "receipt": receipt,
        "signature": signature
    }
\end{lstlisting}

\subsection{Consideraciones de Concurrencia}

\begin{itemize}[leftmargin=*]
    \item \textbf{SELECT FOR UPDATE:} Bloquea filas durante transacción
    \item \textbf{Nivel de aislamiento:} READ COMMITTED o superior
    \item \textbf{Timeouts:} 5 segundos máximo por transacción
    \item \textbf{Retry logic:} Cliente reintenta una vez en caso de deadlock
\end{itemize}

\section{Implementación Criptográfica}

\subsection{Firma Digital de Recibos (Ed25519)}

\subsubsection{Generación de Claves}

\begin{lstlisting}[language=Python]
import nacl.signing
import nacl.encoding

# Generar par de claves (una sola vez, al inicio)
signing_key = nacl.signing.SigningKey.generate()
verify_key = signing_key.verify_key

# Almacenar en variables de entorno
PRIVATE_KEY = signing_key.encode(encoder=nacl.encoding.Base64Encoder)
PUBLIC_KEY = verify_key.encode(encoder=nacl.encoding.Base64Encoder)
\end{lstlisting}

\subsubsection{Firma de Recibo}

\begin{lstlisting}[language=Python]
import json
import nacl.signing
import nacl.encoding

def create_and_sign_receipt(enrollment):
    # Crear estructura del recibo
    receipt = {
        "student_id": str(enrollment.student_user_id),
        "project_id": str(enrollment.project_id),
        "period_id": str(enrollment.period_id),
        "timestamp": enrollment.created_at.isoformat(),
        "nonce": secrets.token_hex(16)
    }
    
    # Serializar de forma canonica
    receipt_json = json.dumps(receipt, sort_keys=True)
    receipt_bytes = receipt_json.encode('utf-8')
    
    # Cargar clave privada
    signing_key = nacl.signing.SigningKey(
        settings.PRIVATE_KEY,
        encoder=nacl.encoding.Base64Encoder
    )
    
    # Firmar
    signed = signing_key.sign(receipt_bytes)
    signature = signed.signature
    
    # Codificar firma en base64
    signature_b64 = base64.b64encode(signature).decode('utf-8')
    
    return receipt, signature_b64
\end{lstlisting}

\subsubsection{Verificación de Recibo}

\begin{lstlisting}[language=Python]
def verify_receipt(receipt, signature_b64):
    try:
        # Reconstruir bytes del recibo
        receipt_json = json.dumps(receipt, sort_keys=True)
        receipt_bytes = receipt_json.encode('utf-8')
        
        # Decodificar firma
        signature = base64.b64decode(signature_b64)
        
        # Cargar clave publica
        verify_key = nacl.signing.VerifyKey(
            settings.PUBLIC_KEY,
            encoder=nacl.encoding.Base64Encoder
        )
        
        # Verificar (lanza excepcion si invalido)
        verify_key.verify(receipt_bytes, signature)
        
        return True
    except Exception as e:
        return False
\end{lstlisting}

\subsection{Hashing de Códigos de Inscripción}

\begin{lstlisting}[language=Python]
import hmac
import hashlib
import secrets
import string

def generate_enrollment_code(length=6):
    """Genera codigo alfanumerico aleatorio"""
    alphabet = string.ascii_uppercase + string.digits
    return ''.join(secrets.choice(alphabet) for _ in range(length))

def hash_code(code: str) -> str:
    """Hashea codigo usando HMAC-SHA256"""
    key = settings.HMAC_SECRET_KEY.encode('utf-8')
    message = code.encode('utf-8')
    h = hmac.new(key, message, hashlib.sha256)
    return h.hexdigest()

def verify_code(code: str, stored_hash: str) -> bool:
    """Verifica codigo contra hash almacenado"""
    computed_hash = hash_code(code)
    return hmac.compare_digest(computed_hash, stored_hash)
\end{lstlisting}

\subsection{Generación de QR de Entrada}

\begin{lstlisting}[language=Python]
import jwt
import qrcode
import io
import base64

def generate_slot_qr(slot_registration):
    # Crear token JWT
    payload = {
        "sub": str(slot_registration.student_user_id),
        "slot_reg_id": str(slot_registration.id),
        "slot_id": str(slot_registration.slot_id),
        "exp": slot_registration.slot.starts_at.timestamp(),
        "type": "slot_checkin"
    }
    
    token = jwt.encode(
        payload,
        settings.JWT_SECRET_KEY,
        algorithm="HS256"
    )
    
    # Generar imagen QR
    qr = qrcode.QRCode(version=1, box_size=10, border=4)
    qr.add_data(token)
    qr.make(fit=True)
    
    img = qr.make_image(fill_color="black", back_color="white")
    
    # Convertir a base64
    buffer = io.BytesIO()
    img.save(buffer, format="PNG")
    img_b64 = base64.b64encode(buffer.getvalue()).decode()
    
    return {
        "qr_token": token,
        "qr_data_url": f"data:image/png;base64,{img_b64}"
    }
\end{lstlisting}

\section{Conclusión}

Este documento de diseño técnico proporciona una guía completa para la implementación del sistema Secure Fair. La arquitectura propuesta es escalable, segura y cumple con los requisitos tanto funcionales como criptográficos. El esquema de base de datos normalizado garantiza integridad referencial, mientras que las transacciones atómicas previenen condiciones de carrera. La implementación criptográfica con Ed25519 y HMAC proporciona las garantías de seguridad necesarias para un sistema de registro crítico.

\end{document}
